% Part: lambda-calculus
% Chapter: introduction
% Section: conversion

\documentclass[../../../include/open-logic-section]{subfiles}

\begin{document}

\olfileid{lam}{int}{con}
\olsection{રૂપાંતરણ અને લઘુકરણ}

લૅમ્બડા કલનમાં રૂપાંતરણ કે લઘુકરણના ત્રણ
પ્રકાર છે. આગળ જોઈશું તેમ, રૂપાંતરણ અને
લઘુકરણ વચ્ચેનો તફાવત એ છે કે
``બદલાવ'' એકદિશ છે કે દ્વિદિશ.

$\alpha$-રૂપાંતરણ ચલનું નામ બદલવા વિશે છે.
$\lambda x. x$ અને $\lambda y.y$ એ એક
જ વિધેય, એટલે ઓળખ-વિધેય, દર્શાવે છે એમ
માનવું સ્વાભાવિક છે. આ વિચારને નીચે મુજબ
ઔપચારિક સ્વરૂપ આપવાનું છે:
\sourcecorrection{OLMD-188}{સ્થિર મૂળ આ
દ્વિબિંદુ પછી તરત સમાપ્ત થાય છે; જણાવેલી
ઔપચારિક વ્યાખ્યા આ વિભાગમાં આપેલી નથી.
અલગ $\alpha$-રૂપાંતરણ વિભાગમાં વિગત છે,
પણ તેને અહીં લખાયેલું ગણ્યું નથી.}

\end{document}
